\documentclass[journal]{IEEEtran}

\usepackage{amsmath,amssymb,bm}
\usepackage{array}
\usepackage{booktabs}
\usepackage{cite}
\usepackage{graphicx}
\usepackage{multirow}
\usepackage{url}
\usepackage[hidelinks]{hyperref}
\usepackage{amsthm}

\newtheoremstyle{ieeestyle}{}{}{\itshape}{}{\itshape}{:}{.5em}{}
\theoremstyle{ieeestyle}
\newtheorem{assumption}{Assumption}
\newtheorem{lemma}{Lemma}
\newtheorem{proposition}{Proposition}
\newtheorem{corollary}{Corollary}
\newtheorem{theorem}{Theorem}

\interdisplaylinepenalty=2500
\graphicspath{{figures/}}

\newcommand{\R}{\mathbb{R}}
\newcommand{\trans}{\mathsf{T}}
\newcommand{\diag}{\operatorname{diag}}
\newcommand{\normtwo}[1]{\left\lVert #1\right\rVert_2}
\newcommand{\norminf}[1]{\left\lVert #1\right\rVert_\infty}

\title{Integrated Lane Keeping and Vehicle Collision Avoidance:
Covariant Ego--Target Estimation and Conditional Control Analysis}

\author{Author Name(s)%
\thanks{Author names, affiliations, corresponding-author information, and
funding acknowledgments must be completed before submission. This manuscript
is a research draft based on the accompanying MATLAB repository.}}

\markboth{Draft Manuscript---September 2026}%
{Author Name(s): Covariant Ego--Target Estimation}

\hypersetup{
    pdftitle={Integrated Lane Keeping and Vehicle Collision Avoidance:
Covariant Ego--Target Estimation and Conditional Control Analysis},
    pdfauthor={Author Name(s)},
    pdfsubject={Integrated lane keeping and vehicle collision avoidance
under bounded state-estimation uncertainty},
    pdfkeywords={autonomous vehicles, lane keeping, collision avoidance,
safety guarantees, state estimation}
}

\begin{document}

\maketitle

\begin{abstract}
This paper develops a measured-input estimator and a predictive controller
for lane keeping with vehicle collision avoidance. The estimator retains
the ego-motion-aware nonlinear relative motion model of the prior
high-gain design but not its differentiated ego coordinates: measured
acceleration and yaw rate remain exogenous inputs, so neither zero ego
jerk nor zero ego yaw acceleration is assumed. A direct body-velocity
measurement follows from GNSS speed and the kinematic single-track
relation under a bounded mismatch. Combining the shared gyro error before
bounding it gives a two-stage velocity--target ISS certificate; a parallel
yaw observer supplies orientation for inertial reconstruction. The target
injection retains its normalized observer LMI, while a free Lyapunov
metric and two Lipschitz multipliers certify decay at a lower bandwidth.
The controller jointly optimizes the input sequence and separation angles
in one convex inner restriction of the held affine bicycle model. A soft
sampled CLF governs tracking, and hard node-separation, actuator, slew,
chart and terminal constraints govern acceptance. The terminal set is a
robust modal cruise set reached by a fixed encounter-exit deadline.
Retaining the complete shifted certificate makes compatible successor
subproblems feasible. Synthetic trials quantify the estimator's reduction
in derivative noise and distinguish controller admission, node safety and
inter-node collision. The estimation theorem is continuous-time; the
control certificate is conditional on the declared sampled plant and does
not establish whole-hold safety or timely admission of every encounter.
\end{abstract}

\begin{IEEEkeywords}
Autonomous vehicles, collision avoidance, cascaded estimator, control
Lyapunov function, high-gain observer, input-to-state stability, nonlinear
relative motion model, predictive control, recursive feasibility,
set-membership estimation, vehicle state estimation.
\end{IEEEkeywords}

\section{Introduction}

\IEEEPARstart{L}{ane}-keeping systems (LKSs) have been extensively studied
and widely deployed in production vehicles \cite{WeiPfefferEdelmann2024}.
Conventional LKSs, however, primarily regulate the ego vehicle's lateral
displacement and heading relative to a lane reference. When an obstacle vehicle
poses a collision risk, avoidance is often left to the driver. This limitation
has motivated considerable research on integrated lane-keeping and vehicle
collision-avoidance systems
\cite{TurriCarvalhoTsengJohanssonBorrelli2013,FahmyGhanyBaumann2018,
ChengLiLiuLiGuo2021,BruggemannSteevesKrstic2022}. These systems retain lane
keeping as the nominal behavior and initiate an automated evasive response when
a collision risk arises.

Collision avoidance has also been studied extensively, but most existing
studies assume that the motion states of the ego and obstacle vehicles are
accurately available. This assumption ignores estimation uncertainty, which can
invalidate controller safety guarantees and increase collision risk. A
safety-critical controller must therefore account for errors in the estimated
states of both vehicles, especially in collision-avoidance scenarios, in which
the ego vehicle often performs complex maneuvers that make the estimation of
surrounding-vehicle states nontrivial. Alai \emph{et al.} developed high-gain nonlinear observers
for obstacle-vehicle trajectory estimation and demonstrated the importance of an
accurate relative motion model for such tracking
\cite{AlaiZemoucheRajamani2024,AlaiSharmaAlexanderRajamani2024}. Sharma
\emph{et al.} subsequently showed that reliable obstacle tracking under complex
ego maneuvers requires a nonlinear relative motion model (NRMM) that explicitly
accounts for ego motion \cite{SharmaAlaiRajamani2026}. Their observer
differentiates that model into third-order companion coordinates, which reduces
the design to a single high-gain theorem but introduces derivatives of the ego
acceleration and yaw rate and therefore assumes that the ego vehicle has zero
body-frame jerk and zero yaw acceleration. These assumptions do not generally
hold in the collision-avoidance scenarios
considered here. This paper therefore develops an estimator that retains their
relative motion model but not its differentiated ego coordinates, so that
neither zero ego jerk nor zero ego yaw acceleration is assumed.

Beyond estimation uncertainty, collision-avoidance performance also depends
on the fidelity of the vehicle model and the generality of the maneuver
formulation. To retain online tractability, many existing controllers rely on
reduced-order models and restricted operating assumptions. Examples include
kinematic bicycle or locally linear lateral models whose validity is tied to
small slip, small lateral acceleration, or low-curvature highway operation
\cite{HeZengZhangSreenath2021,TurriCarvalhoTsengJohanssonBorrelli2013,
BruggemannSteevesKrstic2022}. Other formulations use simplified acceleration
models or two-degree-of-freedom vehicle dynamics under small-angle and linear
tire-force approximations
\cite{ChengLiLiuLiGuo2021,AmmourOrjuelaBasset2022,
DingZhuangHuEtAl2024}. These abstractions are appropriate within their intended
operating envelopes, but evasive steering and braking couple the longitudinal
and lateral vehicle dynamics through tire forces and load transfer
\cite{GuoHuWang2016,FahmyGhanyBaumann2018,HeLiuLvJiLiu2019,
BealBoyd2019}. Moreover, the available tire force becomes strongly nonlinear
as the tires approach saturation
\cite{ParkGerdes2017,ForsAnistratovOlofssonNielsen2021,
ThompsonDallasGohBalachandran2024}.
A further limitation is that collision avoidance is often organized around a
prescribed set of maneuvers. Finite-state or discrete high-level controllers
may select among lane keeping, braking, overtaking, left or right lane changes,
and returning to the original lane
\cite{HeZengZhangSreenath2021,AmmourOrjuelaBasset2022,
ChenAksunGuvenc2025}. Other methods evaluate a predetermined single-lane-change
trajectory or a small number of straight- or low-curvature-road obstacle
configurations
\cite{HeLiuLvJiLiu2019,TurriCarvalhoTsengJohanssonBorrelli2013,
BruggemannSteevesKrstic2022,DingZhuangHuEtAl2024}. Such formulations can be
effective in the scenarios for which they are designed, and dedicated methods
have also been developed for emergency avoidance on curved roads
\cite{LaiYang2023}. However, accommodating oblique cut-ins or interactions that
do not conform to predefined lane and vehicle roles may require additional
maneuver modes, transition logic, or scenario-specific redesign. These
limitations motivate a unified collision-avoidance controller based on a
dynamic vehicle model that continuously coordinates steering and
braking without prescribing the avoidance maneuver in advance.

Model fidelity and maneuver flexibility alone do not establish closed-loop
safety. Several MPC-based collision-avoidance controllers impose geometric or
dynamic constraints over a finite prediction horizon and demonstrate effective
evasive maneuvers, but do not prove that collision-free feasibility is
preserved after each receding-horizon update
\cite{TurriCarvalhoTsengJohanssonBorrelli2013,
ForsAnistratovOlofssonNielsen2021,AmmourOrjuelaBasset2022,
ThompsonDallasGohBalachandran2024}. Control barrier functions (CBFs) provide
sufficient conditions for rendering a prescribed safe set forward invariant
\cite{AmesXuGrizzleTabuada2017}, but this guarantee remains conditional on the
existence of an admissible input satisfying the barrier condition. Conventional
CBF--QP safety filters enforce this condition pointwise and typically minimize
the instantaneous deviation from a nominal input. Consequently, feasibility at
the current state does not ensure that safe inputs remain available at
subsequent states \cite{ZengZhangLiSreenath2021}, and a pointwise filter may
fail to recognize an evasive trajectory that temporarily reduces longitudinal
separation while building sufficient lateral clearance, resulting in myopic or
deadlocked behavior \cite{BlackJankovicSharmaPanagou2023}. Predictive CBF and
MPC--CBF formulations mitigate this limitation by incorporating future system
evolution. Zeng \emph{et al.} imposed discrete-time CBF constraints over an MPC
horizon, but identified persistent feasibility as an open problem for a fixed
CBF decay rate and assumed perfect vehicle-position estimates in their racing
example \cite{ZengZhangSreenath2021}. Breeden and Panagou constructed a
path-based predictive CBF, whereas Wabersich and Zeilinger and Huang
\emph{et al.} derived predictive CBFs from finite-horizon optimal-control value
functions associated with terminal safe sets
\cite{BreedenPanagou2022,WabersichZeilinger2023,
HuangWangMargellosGoulart2025}. The presence of a prediction horizon alone,
however, does not guarantee that a collision-free optimization remains feasible
under repeated receding-horizon updates, particularly when state-estimation
uncertainty is considered.

We carry the safe continuation itself as controller state. The input plan,
occupied-set enclosures and separating directions are retained together,
so a convex majorant touching at the shifted certificate contains a
feasible witness for the next sample. Lane keeping enters through a soft
sampled CLF and a quadratic tracking objective. Collision separation,
actuator limits and terminal membership remain hard. The terminal
condition combines exit from the current sensing region by a fixed
deadline with a robust invariant cruise set, and a current observation
must confirm target release. This construction gives conditional
recursive feasibility after admission; a single convex restriction need
not admit every feasible encounter. Section~\ref{sec:estimator} develops
the estimator and its error enclosure, Section~\ref{sec:control} the
controller that consumes it, and Section~\ref{sec:results} the synthesized
design and recorded trials. The estimation theorem is continuous-time;
the control certificate covers prediction nodes of the declared model.

\section{System Overview}

\begin{figure}[!t]
    \centering
    \includegraphics[width=0.94\columnwidth]
        {driving_environment_curved_v08.pdf}
    \caption{Representative two-lane driving environment. The road-edge
    lines separate the travel lanes from curb-side emergency lanes, each wide
    enough to accommodate a vehicle, while the road centerline lies midway
    between the curb-defined boundaries. The ego vehicle follows the
    nominal path, with an obstacle vehicle cutting into the ego
    vehicle's lane from the adjacent lane.}
    \label{fig:driving_environment}
\end{figure}

We consider an ego vehicle traveling on a road whose drivable surface is
bounded by road curbs. A nominal cruise path is supplied as known route
geometry, while other vehicles travel within
or enter the sensed road environment, as illustrated in
Fig.~\ref{fig:driving_environment}. The reference may be straight, circular,
or a declared smooth curvature profile. The ego vehicle is equipped with a global
navigation satellite system (GNSS) receiver for global position and velocity
measurements and with an inertial measurement unit (IMU) and a gyroscope
for body-frame
acceleration and yaw-rate measurements, respectively. LiDAR provides
geometric measurements of road curbs together with
relative-position measurements of obstacle vehicles, while camera-based
perception detects obstacle vehicles. The problem addressed
in this paper is to design an integrated perception--estimation--control
system such that, in the absence of a predicted collision risk, the ego
vehicle follows the nominal path at a prescribed constant speed.
When continuing this nominal behavior is predicted to create a collision
risk, the system coordinates steering and longitudinal actuation to maintain
collision separation while keeping the ego-vehicle footprint within the
curb-defined drivable surface, and subsequently returns the vehicle to the
nominal path and prescribed speed once the risk has been resolved. Thus, the
curb-defined road boundary is a physical constraint of the intended system,
whereas the nominal path and reference speed are recoverable nominal
objectives. The controller developed below addresses the reference-following
and vehicle-separation part of this problem. Its terminal construction
currently requires a domain without physical road boundaries; the curb
perception module therefore does not yet close that part of the safety
argument.

\begin{figure*}[!t]
    \centering
    \includegraphics[width=0.92\textwidth]
        {ego_obstacle_dynamics_v02.pdf}
    \caption{Ego--obstacle dynamics. The ego vehicle (blue) and an obstacle
    vehicle (orange) travel on courses that may lead to a collision. The
    vehicle frames are centered at the vehicle centers of mass, the yaw
    angles \(\psi_E\) and \(\psi_C\) orient them relative to the inertial
    \(X_I\) axis, the side slip angles \(\beta_E\) and \(\beta_C\)
    separate the velocities \(V_E\) and \(V_C\) from the vehicle
    longitudinal axes, and the course angles \(\gamma_E\) and
    \(\gamma_C\) orient the velocities relative to the inertial \(X_I\)
    axis.}
    \label{fig:ego_obstacle_frames}
\end{figure*}

Fig.~\ref{fig:ego_obstacle_frames} details the vehicle dynamics behind
the scenario of Fig.~\ref{fig:driving_environment}: an ego vehicle
with speed \(V_E\) and yaw angle \(\psi_E\) and an obstacle vehicle
with speed \(V_C\) and yaw angle \(\psi_C\) travel on courses that may
lead to a collision. The controller declares one ego model, the dynamic
bicycle with locally linearized modified Fiala tires; the estimator uses only its kinematics and
the kinematic single-track relation between yaw rate, speed and sideslip,
so no cornering stiffness enters any estimation certificate. With
\(\bm p_E\in\R^2\) and \(\bm v_E\in\R^2\) the inertial position and
velocity,
\(\bm v^E=[v_{x,E}\;\;v_{y,E}]^{\trans}=R(\psi_E)^{\trans}\bm v_E\)
the body velocity, and
\begin{equation}
    R(\psi_E)=\begin{bmatrix}
        \cos\psi_E&-\sin\psi_E\\
        \sin\psi_E&\cos\psi_E
    \end{bmatrix}
    \label{eq:planar_rotation}
\end{equation}
the body-to-inertial rotation, the position kinematics are
\begin{equation}
    \dot{\bm p}_E=\bm v_E=R(\psi_E)\bm v^E .
    \label{eq:position_kinematics}
\end{equation}
The body-frame force and moment balances are
\begin{equation}
    \begin{aligned}
    m_E\dot v_{x,E}&=F_{xf}^{b}+F_{xr}-F_R(v_{x,E})
        +m_Ev_{y,E}\dot\psi_E+m_Eb,\\
    m_E\dot v_{y,E}&=F_{yf}^{b}+F_{yr}-m_Ev_{x,E}\dot\psi_E,\\
    I_{z,E}\ddot\psi_E&=l_{f,E}F_{yf}^{b}-l_{r,E}F_{yr},
    \end{aligned}
    \label{eq:bicycle_model}
\end{equation}
where $[F_{xf}^{b},F_{yf}^{b}]^{\trans}
=R(\delta_F)[F_{xf},F_{yf}]^{\trans}$ rotates the front wheel forces into
the body frame, $F_R$ is passive road load, and $b$ is a declared
longitudinal acceleration bias. The slip convention and modified Fiala
forces are
\begin{equation}
    \begin{aligned}
    \alpha_f&=\operatorname{atan2}(v_{y,E}+l_{f,E}\dot\psi_E,v_\tau)-\delta_F,
    & F_{yf}&=\Phi_f(\alpha_f;Q_f),\\
    \alpha_r&=\operatorname{atan2}(v_{y,E}-l_{r,E}\dot\psi_E,v_\tau),
    & F_{yr}&=\Phi_r(\alpha_r;Q_r).
    \end{aligned}
    \label{eq:tire_forces}
\end{equation}
Here $v_\tau=\max(v_{x,E},v_{\rm f})$ with $v_{\rm f}>0$ regularizes
only the tire denominator. For axle \(i\in\{f,r\}\), let \(t=\tan\alpha\), \(C_i>0\) be its
nominal cornering stiffness, and \(Q_i=\eta_i\mu_iF_{zi}\). Following
\cite{FahmyGhanyBaumann2018}, Section II-A, equations (7), (11) and (13),
\begin{equation}
    \Phi_i(\alpha;Q_i)=
    \begin{cases}
    -C_it+\dfrac{C_i^2}{3Q_i}|t|t-\dfrac{C_i^3}{27Q_i^2}t^3,
       & |\alpha|<\alpha_{{\rm sl},i},\\[3pt]
    -Q_i\operatorname{sgn}(\alpha), & |\alpha|\geq\alpha_{{\rm sl},i},
    \end{cases}
    \label{eq:fiala_force}
\end{equation}
where \(\alpha_{{\rm sl},i}=\arctan(3Q_i/C_i)\),
\(\eta_i=\sqrt{1-\beta^2}\), and \(F_{xi}=\beta\mu_iF_{zi}\)
for \(|\beta|<1\). The same signed input \(\beta\) applies to both axles: negative values brake and positive values request throttle. Static loads are
\(F_{zf}=m_Egl_{r,E}/L\) and \(F_{zr}=m_Egl_{f,E}/L\), with
\(L=l_{f,E}+l_{r,E}\). Each axle represents an equivalent wheel pair.
Here \(m_E\) and \(I_{z,E}\) denote mass and yaw inertia; the lever arms
are \(l_{f,E},l_{r,E}\), and \(\delta_F\) is front steering.
We use the symmetric branch condition and the cited paper's equation (13);
the additional arctangent printed in its equation (12) is inconsistent
with (13). At \(|\beta|=1\), the force law is defined by
\(Q_i=0\), \(\Phi_i=0\), without evaluating its singular polynomial.
Joint derivatives are singular at these endpoints; linearization anchors
are kept strictly inside the input interval.
The controller linearizes this force locally as specified in
Section~\ref{subsec:frenetModel}; separate axle friction-limit analysis
is outside the present controller formulation.
The speed, sideslip, and course of
Fig.~\ref{fig:ego_obstacle_frames} are
\begin{equation}
    V_E=\normtwo{\bm v_E},
    \qquad
    \beta_E=\arctan\frac{v_{y,E}}{v_{x,E}},
    \qquad
    \gamma_E=\psi_E+\beta_E .
    \label{eq:yaw_from_course}
\end{equation}
The obstacle vehicle has its own unobserved inputs. The tracker uses the
relative kinematics of Section~\ref{subsec:targetTracker}; the controller
propagates its admitted bounded-motion set independently of the ego tire model.

\begin{figure*}[!t]
    \centering
    \includegraphics[width=0.96\textwidth]{system_architecture.pdf}
    \caption{Information flow in the proposed
    perception--estimation--control architecture. Camera and LiDAR data are
    jointly processed for vehicle detection, while LiDAR data also supply
    road-curb extraction. The estimator combines the
    resulting obstacle-relative positions with GNSS, IMU, and gyroscope signals.
    The estimator produces separate ego-vehicle and
    obstacle-vehicle state estimates. The ego estimate is supplied to the
    MPC, while the obstacle-vehicle estimate enters the safety
    constraint. The nominal cruise path, taken from an existing route
    map, and road boundary define the nominal objective and safety constraint,
    respectively, used by the predictive controller to compute steering and
    longitudinal actuation commands. The road-boundary branch shows the
    intended interface; the current terminal certificate requires a domain
    without physical road boundaries.}
    \label{fig:architecture}
\end{figure*}

The architecture that realizes this behavior comprises perception,
estimation, and control modules, as summarized in Fig.~\ref{fig:architecture}.
The perception module processes synchronized camera and LiDAR data jointly
for obstacle-vehicle detection and relative localization. LiDAR data are also
passed to a dedicated road-curb extraction component, whose output specifies
the curb-defined road boundary; the nominal path is known route geometry and
requires no extraction. The estimator combines the obstacle-relative
positions with the GNSS, IMU, and gyroscope measurements to estimate the motion
states of the ego vehicle and each tracked obstacle vehicle. Each
estimate is accompanied by a componentwise error enclosure, valid under
the declared sensor, model and operating-domain premises, with which the
controller tightens its own constraints. The enclosure bounds the current
estimate; propagating it over the prediction horizon is the controller's
own obligation. The road-boundary branch represents the intended perception interface;
physical boundary constraints remain outside the current terminal
certificate. The known nominal path supplies the control reference. At each
sampling instant, the
controller constructs the nominal objective and safety constraints from the
reference geometry, estimated motion states, their enclosures, and the
prescribed reference speed, solves one convex program, checks its result
independently, and issues steering and longitudinal actuation commands to
the ego vehicle.

\section{Perception}
\label{sec:perception}

The perception module obtains, from the onboard cameras and LiDAR, the
positions of the obstacle vehicles and of the road curbs relative to the ego
vehicle.

\begin{figure*}[!t]
    \centering
    \begin{minipage}[t]{0.605\textwidth}
        \centering
        \includegraphics[width=\textwidth]{yolo_lidar_projection_v01.png}\\[2pt]
        (a)
    \end{minipage}\hfill
    \begin{minipage}[t]{0.34\textwidth}
        \centering
        \includegraphics[width=\textwidth]{obstacle_rectangle_fit_topdown_v01.png}\\[2pt]
        (b)
    \end{minipage}
    \caption{Visualization of the perception outputs. (a) The LiDAR point
    cloud projected into the scene. The points associated with the YOLO
    detection of the obstacle vehicle are shown in yellow, the curb
    points detected in this frame in red, and the tracked curb
    quadratics as orange curves, while the green dashed line is the
    nominal route. The inset shows the front-camera image with the
    detected bounding box. (b) Top view
    of the points reflected from the obstacle vehicle and the rectangle
    fitted to their outermost points, whose center marks the measured
    relative position.}
    \label{fig:obstacle_perception}
\end{figure*}

\subsection{Obstacle-Vehicle Perception}
Obstacle vehicles are detected in the camera images by a YOLO
detector \cite{RedmonDivvalaGirshickFarhadi2016}, which returns a
two-dimensional bounding box for each of them. The LiDAR point cloud
is then projected onto the image, as shown in
Fig.~\ref{fig:obstacle_perception}(a). Each LiDAR point \(\bm q\) is
transformed into the camera frame and from there into pixel
coordinates,
\begin{equation}
    \begin{gathered}
    \begin{bmatrix}X&Y&Z\end{bmatrix}^{\trans}
    =\bm R_{CL}\,\bm q+\bm t_{CL},\\
    \begin{bmatrix}u\\v\end{bmatrix}
    =\frac{f}{Z}
    \begin{bmatrix}X\\Y\end{bmatrix}
    +\begin{bmatrix}W/2\\H/2\end{bmatrix},
    \end{gathered}
    \label{eq:lidar_camera_projection}
\end{equation}
where \((\bm R_{CL},\bm t_{CL})\) are the LiDAR-to-camera extrinsics,
\((X,Y,Z)\) are the point coordinates in the camera frame with \(Z>0\)
the optical depth, \(f\) is the focal length, and \(W\) and
\(H\) are the image width and height. The LiDAR points whose pixel coordinates \((u,v)\) fall
inside a detected bounding box are associated with that vehicle, and
the ground and other extraneous points among them are removed,
leaving the points reflected from the obstacle vehicle.

In this paper, the physical outline of a vehicle is regarded as a
rectangle in the top view. A rectangle is therefore fitted to the outermost
of the retained points, as shown in Fig.~\ref{fig:obstacle_perception}(b), and
the center of the fitted rectangle yields the measured relative
position of the obstacle vehicle,
\begin{equation}
    \bm\rho=\begin{bmatrix}\rho_x&\rho_y\end{bmatrix}^{\trans}\in\R^2 ,
    \label{eq:perception_relative_position}
\end{equation}
expressed in the ego-vehicle frame.

\subsection{Road-Curb Extraction and Tracking}

\subsubsection{Local Geometric Feature Analysis}
A curb is a step of a few centimetres in an otherwise smooth road
surface, so the local geometry of the LiDAR returns changes abruptly
where the beams cross it. Four local geometric features are considered
to capture this discontinuity---adjacent-ring compression, vertical
step, radial jump, and local shape. In the following, \(\bm p_{i,j}\)
and \(z_{i,j}\) denote the horizontal position and height of the
return at elevation ring \(i\) and azimuth \(j\), and
\(d_{i,j}=\|\bm p_{i,j}\|\) its horizontal range.

\paragraph{Adjacent-ring compression}
Flat ground predicts the horizontal spacing between consecutive rings
from the sensor height and beam angles alone, and a raised curb face
compresses the spacing \(\Delta_{i,j}\) observed between ring \(i\)
and its upper neighbor below that prediction \(\hat\Delta_{i}\),
\begin{equation}
    \gamma_{i,j}
    =\frac{\hat\Delta_{i}-\Delta_{i,j}}{\hat\Delta_{i}} .
    \label{eq:curb_ring_compression}
\end{equation}

\paragraph{Vertical step}
The height of the curb face appears directly as a step between
same-ring neighbors,
\begin{equation}
    \sigma_{i,j}
    =\max\bigl(|z_{i,j}-z_{i,j-1}|,\,|z_{i,j+1}-z_{i,j}|\bigr).
    \label{eq:curb_vertical_step}
\end{equation}

\paragraph{Radial jump}
On flat ground the horizontal range varies smoothly with azimuth,
whereas a beam climbing the curb face is pulled inward, which breaks
the smoothness captured by the second difference
\begin{equation}
    \zeta_{i,j}
    =|d_{i,j-1}-2d_{i,j}+d_{i,j+1}| .
    \label{eq:curb_radial_jump}
\end{equation}

\paragraph{Local shape}
On the road surface a ring sweeps an arc whose chord is perpendicular
to the line of sight, whereas along a curb the returns align with the
road edge, so the normalized projection of the chord of the same-ring
neighbors onto the line of sight
\begin{equation}
    \tau_{i,j}
    =\frac{\bigl|\bm p_{i,j}^{\trans}
        (\bm p_{i,j+1}-\bm p_{i,j-1})\bigr|}
      {\|\bm p_{i,j}\|\,\|\bm p_{i,j+1}-\bm p_{i,j-1}\|}
    \label{eq:curb_local_shape}
\end{equation}
rises where the ring turns along the curb.

\begin{figure}[!t]
    \centering
    \includegraphics[width=\columnwidth]{curb_ring_features_v01.pdf}
    \caption{The four local geometric features along the forward half of
    ring 43 of a 64-ring LiDAR mounted \(2.35\) m above a \(7\) m road
    bounded by \(0.15\) m curbs, with the ego vehicle on the road
    centerline. The ring meets flat ground at a radius of \(7.85\) m. The
    dashed lines mark the azimuth bins at which it reaches the curbs,
    predicted from that radius and the road width alone.}
    \label{fig:curb_ring_features}
\end{figure}

Fig.~\ref{fig:curb_ring_features} follows the four features along one
ring. Each holds a flat baseline over the road
surface and over the raised walkway, and all four rise by more than an
order of magnitude at the same places, where the ring crosses
the curbs.

\subsubsection{Curb Fitting}
The road heading \(\varphi\) is estimated from the single longest
near-linear structure among the candidates and is held near a prior
propagated with the ego motion. It defines a road-aligned frame in
which each curb is the quadratic
\begin{equation}
    y_{\mathrm{b}}(x)=c_{2}x^{2}+c_{1}x+c_{0},
    \qquad x\in[x_{\min},x_{\max}],
    \label{eq:curb_quadratic}
\end{equation}
one curve per side over a finite forward domain. The coefficients are
fitted to the side's candidate points by least squares.

\subsubsection{Curb Tracking}
A single-frame fit can latch onto a structure that merely resembles a curb and jump far from the tracked
geometry, so the fit of frame \(k\), with coefficients \(c_{i,k}\), is
admitted only through an amplitude-limiting filter on the coefficient
increments,
\begin{equation}
    |c_{i,k}-c_{i,k-1}|\le T_i,
    \qquad i\in\{0,1,2\},
    \label{eq:curb_amplitude_limit}
\end{equation}
and a fit violating any bound leaves the frame to the prediction. The
curve itself is tracked through three points \(\bm p_s\) placed on it
at constant longitudinal stations \(x_s\). Between frames each point
is propagated by the measured ego motion,
\begin{equation}
    \bm p_s^-=A\,\bm p_s+\bm b,
    \qquad
    P_s^-=A P_s A^{\trans}+Q,
    \label{eq:curb_point_prediction}
\end{equation}
where \((A,\bm b)\) is the rigid transform from the preceding to the
current road frame, \(P_s\) is the point covariance, and \(Q\) is the
process noise. The admitted fit, evaluated at the same stations,
yields the measured points \(\bm z_s\) with noise covariance \(R\),
and each tracked point follows the Kalman update
\begin{equation}
    K_s=P_s^-\,(P_s^-+R)^{-1},
    \qquad
    \bm p_s=\bm p_s^-+K_s(\bm z_s-\bm p_s^-).
    \label{eq:curb_point_update}
\end{equation}

\section{Estimator Design}
\label{sec:estimator}

The ego body-velocity observer is designed first and the target tracker
second. Their dynamics form a two-stage cascade; the parallel yaw observer
and inertial position output are used for reconstruction and do not feed
this cascade. The measurements are
\begin{equation}
    \begin{aligned}
    \bm y_p&=\bm p_E+\bm n_p, &\quad
    \bm v_m&=\bm v_E+\bm n_v,\\
    \bm a_m&=\bm a_E+\bm n_a, &\quad
    u_3&=\dot\psi_E+d_\omega,\\
    \bm y_R&=\bm\rho+\bm n_R,
    \end{aligned}
    \label{eq:ego_measurements}
\end{equation}
with \(\bm\rho\) the ego-frame relative position \eqref{eq:target_coordinates}
and uniform bounds \(\normtwo{\bm n_p}\leq\bar n_p\),
\(\normtwo{\bm n_v}\leq\bar n_v\), \(\normtwo{\bm n_a}\leq\bar n_a\),
\(|d_\omega|\leq\bar n_g\), \(\normtwo{\bm n_R}\leq\bar n_R\). Inertial
biases are compensated upstream, so no observer carries a bias state.

Two constructions of \cite{SharmaAlaiRajamani2026} are replaced. Their
translational observer differentiates the ego kinematics into third-order
companion chains, which requires \(\dot{\bm a}_E\) and \(\ddot\psi_E\) and
therefore assumes both vanish; here \(\bm a_m\) and \(u_3\) remain
exogenous inputs of every stage and nothing is differentiated. Their yaw
observer reads the course as the yaw, which
\(\gamma_E=\psi_E+\beta_E\) makes exact only at \(\beta_E=0\); here the
same relation supplies a direct body-velocity measurement and, when
informative, a certified heading correspondence for the parallel yaw observer.

\begin{assumption}
\label{asm:ego}
\(0\leq V_E|_{\min}\leq V_E\leq V_E|_{\max}\),
$v_{x,E}\geq V_E\cos\beta_{\rm dom}$, and
\(|\dot\psi_E|\leq\bar\omega_E\).
\end{assumption}

\begin{assumption}
\label{asm:singleTrack}
The ego kinematic single-track relation holds with bounded mismatch,
\begin{equation}
    \dot\psi_E=\frac{V_E\sin\beta_E}{l_{r,E}}+d_{\rm st},
    \qquad|d_{\rm st}|\leq\bar d_{\rm st},
    \label{eq:single_track_mismatch}
\end{equation}
and \(|\beta_E|\leq\beta_{\rm dom}<\pi/2\).
\end{assumption}

Assumption~\ref{asm:singleTrack} bounds unmodeled rear-axle slip. The
identification it replaces is its special case \(\beta_E\equiv0\).

\subsection{Certified Course--Yaw Correspondence}
\label{subsec:courseCertificate}

\begin{lemma}
\label{lem:course}
Let \(V_m=\normtwo{\bm v_m}>\bar n_v\), put
\(\underline V=V_m-\bar n_v\), and suppose \(|l_{r,E}u_3|\leq V_m\).
Define
\begin{equation}
    \hat\beta=\arcsin\frac{l_{r,E}u_3}{V_m},
    \qquad
    y_F=\angle\bm v_m-\hat\beta .
    \label{eq:course_definition}
\end{equation}
Then, with
\begin{equation}
    \bar e_{\sin}
    =\frac{l_{r,E}}{\underline V}
    \Bigl(\bar n_g+\bar d_{\rm st}+\frac{|u_3|\bar n_v}{V_m}\Bigr),
    \label{eq:sine_error_bound}
\end{equation}
\begin{equation}
    \bar e_\beta=\max_{\pm}\Bigl|
    \arcsin\bigl(\Pi\bigl[\tfrac{l_{r,E}u_3}{V_m}\pm\bar e_{\sin}\bigr]\bigr)
    -\hat\beta\Bigr|,
    \label{eq:sideslip_error}
\end{equation}
where \(\Pi\) clips to \([-\sin\beta_{\rm dom},\sin\beta_{\rm dom}]\), one has
\(|\beta_E-\hat\beta|\leq\bar e_\beta\) and
\begin{equation}
    |y_F-\psi_E|\leq B_F:=\arcsin\frac{\bar n_v}{V_m}+\bar e_\beta .
    \label{eq:course_certificate}
\end{equation}
\end{lemma}

\begin{IEEEproof}
By \eqref{eq:single_track_mismatch},
\(\sin\beta_E=l_{r,E}(u_3-d_\omega-d_{\rm st})/V_E\), so
\begin{equation}
    \Bigl|\sin\beta_E-\frac{l_{r,E}u_3}{V_m}\Bigr|
    \leq l_{r,E}\Bigl[\frac{|d_\omega|+|d_{\rm st}|}{V_E}
    +|u_3|\frac{|V_m-V_E|}{V_EV_m}\Bigr],
    \label{eq:sine_derivation}
\end{equation}
and \(|V_m-V_E|\leq\bar n_v\), \(V_E\geq\underline V\) give
\eqref{eq:sine_error_bound}. Since \(|\sin\beta_E|\leq\sin\beta_{\rm dom}\),
the true sine lies in the clipped interval, and \(\arcsin\) is increasing
on it, which gives \eqref{eq:sideslip_error}. For the course, \(\bm v_E\)
lies in the disk of radius \(\bar n_v\) about \(\bm v_m\); the disk
excludes the origin, so \(\bm v_E\) lies in its tangent cone and
\(|\angle\bm v_m-\gamma_E|\leq\arcsin(\bar n_v/V_m)\). Subtracting
\(\gamma_E=\psi_E+\beta_E\) yields \eqref{eq:course_certificate}.
\end{IEEEproof}

The correspondence is used only where these informativeness conditions
hold. Otherwise it supplies the whole circle as a yaw enclosure. Neither
case changes the body-velocity innovation developed next, which remains
defined at standstill.

\subsection{Ego Observers}
\label{subsec:egoObserver}

With $J=\left[\begin{smallmatrix}0&-1\\1&0\end{smallmatrix}\right]$,
$\nabla_u\bm z=\dot{\bm z}+uJ\bm z$, and
$\bm v^E=R(\psi_E)^{\trans}\bm v_E$, define
\begin{equation}
    \begin{gathered}
    \phi(V,w)=\sqrt{\max\{V^2-w^2,V^2\cos^2\beta_{\rm dom}\}},\\
    \bm y_v=[\phi(V_m,l_{r,E}u_3),\ l_{r,E}u_3]^{\trans},\\
    \nabla_{u_3}\hat{\bm v}^E=\bm a_m+k_v(\bm y_v-\hat{\bm v}^E).
    \end{gathered}
    \label{eq:ego_observer}
\end{equation}
On the true forward cone, $\phi(V_E,v_{y,E})=v_{x,E}$. The extension
remains defined for every measured speed, including zero, and its lateral
component is not clipped. Thus
$y_{v,2}-v_{y,E}=l_{r,E}(d_\omega+d_{\rm st})$.

\begin{lemma}
\label{lem:egoRows}
Let $\bm e_v=\hat{\bm v}^E-\bm v^E$, $k_v>0$, and
\begin{equation}
    \begin{split}
    C_\omega(k)=\max_{V\in\{V_E|_{\min},V_E|_{\max}\}}
    \Bigl[(V\sin\beta_{\rm dom}+kl_{r,E}\tan\beta_{\rm dom})^2\\
        +(kl_{r,E}-V\cos\beta_{\rm dom})^2\Bigr]^{1/2}.
    \end{split}
    \label{eq:common_gyro_bound}
\end{equation}
Under Assumptions~\ref{asm:ego}--\ref{asm:singleTrack},
\begin{equation}
    \begin{gathered}
    D^+\normtwo{\bm e_v}\leq-k_v\normtwo{\bm e_v}+\bar\delta_v,\\
    \bar\delta_v=\bar n_a+k_v\sec\beta_{\rm dom}\,\bar n_v
      +C_\omega(k_v)\bar n_g
      +k_vl_{r,E}\sec\beta_{\rm dom}\,\bar d_{\rm st}.
    \end{gathered}
    \label{eq:ego_rows}
\end{equation}
\end{lemma}

\begin{IEEEproof}
Put $\nu=V_m-V_E$, so $|\nu|\leq\bar n_v$. Integrating the
piecewise derivatives of $\phi$ along the segment between
$(V_E,v_{y,E})$ and $(V_m,l_{r,E}u_3)$ gives
\begin{equation}
    \begin{gathered}
    \bm y_v-\bm v^E=\alpha\bm e_1\nu
        +l_{r,E}\bm g(\zeta)(d_\omega+d_{\rm st}),\\
    \bm g(\zeta)=[-\zeta,1]^{\trans},\\
    \cos\beta_{\rm dom}\leq\alpha\leq\sec\beta_{\rm dom},\qquad
    |\zeta|\leq\tan\beta_{\rm dom}.
    \end{gathered}
    \label{eq:velocity_measurement_error}
\end{equation}
Continuity covers branch crossings and the origin. Substituting before
taking norms yields
\begin{equation}
    \begin{split}
    \dot{\bm e}_v={}&(-k_vI-u_3J)\bm e_v+\bm n_a+k_v\alpha\bm e_1\nu\\
        &+(k_vl_{r,E}\bm g(\zeta)-J\bm v^E)d_\omega
          +k_vl_{r,E}\bm g(\zeta)d_{\rm st}.
    \end{split}
    \label{eq:velocity_error_derivation}
\end{equation}
The common gyro column has components
$[v_{y,E}-k_vl_{r,E}\zeta,\ k_vl_{r,E}-v_{x,E}]^{\trans}$.
Maximizing its norm over $\zeta$ and the forward cone, then using
convexity of its squared bound in speed, gives
\eqref{eq:common_gyro_bound}. The rotation is skew and contributes zero
to the norm derivative. The remaining terms give \eqref{eq:ego_rows}.
\end{IEEEproof}

The same gyro error enters measurement and propagation. Bounding those
terms separately would discard their cancellation; no statistical
independence is used in Lemma~\ref{lem:egoRows}. In particular, yaw error
does not enter the velocity comparison.

Yaw and inertial position are estimated in parallel:
\begin{equation}
    \begin{aligned}
    \dot{\hat\psi}_E&=u_3+\chi_m k_\psi
        \operatorname{wrap}(y_F-\hat\psi_E),\\
    \dot{\hat{\bm p}}_E&=\bm v_m+k_p(\bm y_p-\hat{\bm p}_E),
    \end{aligned}
    \label{eq:parallel_outputs}
\end{equation}
where $\chi_m=1$ only when the course correspondence is informative.
On a compatible innovation chart the yaw error obeys its usual
gyro-plus-heading comparison; during uninformative intervals only gyro
propagation is available. An independent orientation set $\Psi$ is
propagated with the gyro-error integral and intersected with each
informative circular arc from Lemma~\ref{lem:course}. Disconnected
components are retained, and the published bound is
$B_\psi=\sup_{\theta\in\Psi}d_{S^1}(\hat\psi_E,\theta)$ about the
actual observer estimate. An empty intersection invalidates orientation
outputs; it does not reset the observer. The position error
\(\bm e_p=\hat{\bm p}_E-\bm p_E\) satisfies
\begin{equation}
    \begin{gathered}
    D^+\normtwo{\bm e_p}\leq-k_p\normtwo{\bm e_p}+\bar n_v+k_p\bar n_p,\\
    \limsup_{t\to\infty}\normtwo{\bm e_p}\leq\bar n_p+\bar n_v/k_p.
    \end{gathered}
    \label{eq:position_error_bound}
\end{equation}

\subsection{Covariant Target Tracker}
\label{subsec:targetTracker}

The target model of \cite{SharmaAlaiRajamani2026}---constant scalar
acceleration and constant sideslip, hence constant curvature
\(\kappa=\sin\beta_C/l_{r,C}\)---is retained in covariant companion
coordinates
\begin{equation}
    \bm\rho=R_E^{\trans}(\bm p_C-\bm p_E),\quad
    \bm q=R_E^{\trans}\bm v_C,\quad
    \bm s=R_E^{\trans}\bm a_C,
    \label{eq:target_coordinates}
\end{equation}
in which the exact dynamics are the third-order chain
\begin{equation}
    \begin{gathered}
    \nabla_{\omega_E}\bm\rho=\bm q-\bm v^E,\quad
    \nabla_{\omega_E}\bm q=\bm s,\quad
    \nabla_{\omega_E}\bm s=\Phi(\bm q,\bm s),\\
    \Phi=-\Omega^2\bm q+3A\Omega J\frac{\bm q}{V},\\
    V=\normtwo{\bm q},\quad
    A=\frac{\bm q^{\trans}\bm s}{V},\quad
    \Omega=\frac{(J\bm q)^{\trans}\bm s}{V^2}.
    \end{gathered}
    \label{eq:target_chain}
\end{equation}
Only \(\bm v^E\) and \(\omega_E\) enter, and the only extra terms are those
of the rotating frame; \(A\) and \(\Omega\) are reconstructions, not states.
The declared operating set is
\begin{equation}
    \begin{gathered}
    a\leq\normtwo{\bm q}\leq b,\quad
    |A|\leq A_{\max},\quad
    \normtwo{\bm\rho}\leq\rho_{\max},\\
    |\Omega|\leq\Omega_{\max}=\frac{b\sin\beta_{\max}}{l_{r,C}},\quad
    \normtwo{\bm s}\leq S=\sqrt{A_{\max}^2+b^2\Omega_{\max}^2}.
    \end{gathered}
    \label{eq:target_domain}
\end{equation}

Estimates may peak outside \eqref{eq:target_domain}, so the observer
evaluates the extension \(\Phi_e\) obtained by replacing \(\bm q,\bm s\) in
\(\Phi\) by their ball projections onto radii \(b,S\), clipping \(A\) and
\(\Omega\) to their bounds, and using \(\max(\normtwo{\bm q},a)\) in every
denominator. It agrees with \(\Phi\) on \eqref{eq:target_domain} and is
defined on \(\R^4\).

\begin{lemma}
\label{lem:lipschitz}
With
\(c_A=S/a\), \(c_\Omega=S\max(a^{-2},2b^{-2})\),
\(c_{Q\Omega}=S\max(a^{-1},2b^{-1})\),
\begin{equation}
    \begin{gathered}
    L_q=\Omega_{\max}^2+2\Omega_{\max}c_{Q\Omega}
        +3\Omega_{\max}\Bigl(\frac{A_{\max}}a+c_A\Bigr)\\
        +3A_{\max}c_\Omega,\\
    L_s=\left\lVert\begin{bmatrix}
        0&2\Omega_{\max}\\3\Omega_{\max}&3A_{\max}/a
        \end{bmatrix}\right\rVert_2,
    \end{gathered}
    \label{eq:lipschitz_constants}
\end{equation}
\(\Phi_e\) satisfies, for all \((\bm q,\bm s),(\hat{\bm q},\hat{\bm s})\),
\begin{equation}
    \begin{split}
    \normtwo{\Phi_e(\bm q,\bm s)-\Phi_e(\hat{\bm q},\hat{\bm s})}\\
    \leq L_q\normtwo{\bm q-\hat{\bm q}}+L_s\normtwo{\bm s-\hat{\bm s}} .
    \end{split}
    \label{eq:global_lipschitz}
\end{equation}
\end{lemma}

\begin{IEEEproof}
Projections and clips are nonexpansive, so they only decrease derivative
bounds, and \(\Phi_e\) is continuous and piecewise differentiable; it
suffices to bound its almost-everywhere Jacobians and integrate along the
segment joining the two arguments. Differentiating in the radial and
tangential basis at fixed \(\bm s\), with \(r=\normtwo{\bm q}\), gives
\(|D_q\alpha|\leq S/\max(r,a)\) and \(|D_q\nu|\leq S/\max(r,a)^2\) for
\(r\leq b\) and \(|D_q\alpha|\leq bS/r^2\),
\(|D_q\nu|\leq2bS/r^3\) for \(r\geq b\), whose suprema over \(r>0\) are
\(c_A\) and \(c_\Omega\), while \(\normtwo{\bm Q}|D_q\nu|\leq c_{Q\Omega}\).
Collecting the four terms of \(\Phi_e\) yields \(L_q\). At fixed
\(\bm q\), the entrywise absolute values of the acceleration Jacobian in the
basis aligned with \(\bm q\) are dominated by the displayed matrix, and its
spectral norm dominates the induced norm; the case \(\bm q=0\) follows by
continuity.
\end{IEEEproof}

The tracker is the covariant high-gain chain
\begin{equation}
    \begin{aligned}
    \nabla_{u_3}\hat{\bm\rho}&=\hat{\bm q}-\hat{\bm v}^E
        +l_1\omega\,(\bm y_R-\hat{\bm\rho}),\\
    \nabla_{u_3}\hat{\bm q}&=\hat{\bm s}
        +l_2\omega^2(\bm y_R-\hat{\bm\rho}),\\
    \nabla_{u_3}\hat{\bm s}&=\Phi_e(\hat{\bm q},\hat{\bm s})
        +l_3\omega^3(\bm y_R-\hat{\bm\rho}).
    \end{aligned}
    \label{eq:target_observer}
\end{equation}
Let \(A_c\) be the integrator-chain matrix, \(C=\bm e_1^{\trans}\), and
\(A_l=A_c-\bm lC\). The normalized direction \(\bm l\) solves the observer
LMI in \((\tilde P,\bm Y=\tilde P\bm l,\gamma)\),
\begin{equation}
    \begin{gathered}
    \min\gamma
    \quad\text{s.t.}\quad
    \tilde P\succeq I,\quad
    \Lambda+\Lambda^{\trans}+2\alpha\tilde P\prec0,\\
    \begin{bmatrix}
    \Lambda+\Lambda^{\trans}+I&-\bm Y\\-\bm Y^{\trans}&-\gamma
    \end{bmatrix}\prec0,
    \qquad \Lambda=\tilde PA_c-\bm YC,
    \end{gathered}
    \label{eq:target_lmi}
\end{equation}
an $\alpha$-stability region with a bounded-real radar-to-state gain.
The metric $P$ used to certify the nonlinear observer is chosen separately.
For the estimate-minus-truth target error $\bm e_T$, set
\begin{equation}
    \bm\epsilon=\diag(I_2,\omega^{-1}I_2,\omega^{-2}I_2)\bm e_T,
    \qquad
    W_T=\sqrt{\bm\epsilon^{\trans}(P\otimes I_2)\bm\epsilon}.
    \label{eq:target_metric}
\end{equation}

\begin{lemma}
\label{lem:dissipation}
Let $P\succ0$, \(\bm b_P=P\bm e_3\), \(t_q,t_s>0\), and
\begin{equation}
    \begin{aligned}
    M={}&-\omega(A_l^{\trans}P+PA_l)
       -t_q\frac{L_q^2}{\omega^2}\bm e_2\bm e_2^{\trans}\\
       &-t_sL_s^2\bm e_3\bm e_3^{\trans}
       -(t_q^{-1}+t_s^{-1})\bm b_P\bm b_P^{\trans}.
    \end{aligned}
    \label{eq:dissipation_matrix}
\end{equation}
If \(M\succeq2\lambda_TP\) then \(\dot W_T\leq-\lambda_TW_T\) in the
absence of exogenous error inputs.
\end{lemma}

\begin{IEEEproof}
In the coordinate \(\bm\epsilon\) the injection block becomes
\(\omega A_l\otimes I_2\), and the common rotation \(-u_3(I_3\otimes J)\)
is skew and contributes zero to
\(\tfrac{d}{dt}\bm\epsilon^{\trans}(P\otimes I_2)\bm\epsilon\). The injection contributes
$\omega\bm\epsilon^{\trans}[(A_l^{\trans}P+PA_l)\otimes I_2]\bm\epsilon$.
Splitting $\Delta\Phi/\omega^2$ into its velocity and acceleration
increments, \eqref{eq:global_lipschitz} gives respective bounds
$(L_q/\omega)\normtwo{\bm\epsilon_2}$ and
$L_s\normtwo{\bm\epsilon_3}$. Each enters the derivative through
$2\bm\epsilon^{\trans}(\bm b_P\otimes I_2)$. Applying
\(2\bm z^{\trans}\bm d\leq\normtwo{\bm z}^2/t+t\normtwo{\bm d}^2\) once per
channel gives
\(\tfrac{d}{dt}W_T^2\leq-\bm\epsilon^{\trans}(M\otimes I_2)\bm\epsilon
\leq-2\lambda_TW_T^2\).
\end{IEEEproof}

\subsection{Cascade ISS}
\label{subsec:cascade}

All core couplings are feedforward. With
$\bm h=[\rho_{\max},\ b/\omega,\ S/\omega^2]^{\trans}$,
Cauchy--Schwarz in the metric \eqref{eq:target_metric} gives
\begin{equation}
    \begin{gathered}
    g_{Tv}=\sqrt{P_{11}},\qquad
    g_R=\omega\sqrt{\bm l^{\trans}P\bm l},\\
    g_\omega=\sqrt{\bm h^{\trans}|P|\bm h},\qquad
    g_\Phi=\sqrt{P_{33}}/\omega^2,
    \end{gathered}
    \label{eq:target_disturbance_coefficients}
\end{equation}
where $|P|$ is entrywise. Declared scalar-acceleration and curvature rates
add a forcing bounded by
$\bar\nu_\Phi=(\bar{\dot A}^2+b^4\bar{\dot\kappa}^2)^{1/2}$ to the last
chain equation. Lemmas~\ref{lem:egoRows} and \ref{lem:dissipation} give
\begin{equation}
    \begin{gathered}
    \bm\chi=[\normtwo{\bm e_v},\ W_T]^{\trans},\qquad
    D^+\bm\chi\leq H\bm\chi+\bm d,\\
    H=\begin{bmatrix}-k_v&0\\g_{Tv}&-\lambda_T\end{bmatrix},\qquad
    \bm d=\begin{bmatrix}\bar\delta_v\\
    g_\omega\bar n_g+g_R\bar n_R+g_\Phi\bar\nu_\Phi\end{bmatrix}.
    \end{gathered}
    \label{eq:cascade_comparison}
\end{equation}

\begin{theorem}[Cascade ISS]
\label{thm:cascadeISS}
Let Assumptions~\ref{asm:ego}--\ref{asm:singleTrack}, the hypotheses of
Lemmas~\ref{lem:egoRows}, \ref{lem:lipschitz} and \ref{lem:dissipation},
the true target domain \eqref{eq:target_domain} and the declared continuous
sensor and model-error bounds hold, with $k_v,\lambda_T>0$. Then $H$ is
Metzler and Hurwitz, and $\bm c=-H^{-\trans}\bm 1>0$ gives the linear
copositive ISS function $V_c=\bm c^{\trans}\bm\chi$. For any
$\bm z_0\geq\bm\chi(0)$,
\begin{equation}
    \bm\chi(t)\leq\bm z(t)
    =e^{Ht}\bm z_0+\int_0^t e^{H(t-\tau)}\bm d\,d\tau,
    \label{eq:comparison_solution}
\end{equation}
and $\limsup_{t\to\infty}\bm\chi(t)\leq-H^{-1}\bm d$. The physical
component bounds are
\begin{equation}
    \normtwo{\bm e_{T,i}}\leq c_iW_T,\qquad
    \bm c_T=\diag(1,\omega,\omega^2)\sqrt{\diag(P^{-1})}.
    \label{eq:component_extraction}
\end{equation}
\end{theorem}

\begin{IEEEproof}
$H$ is lower triangular with negative diagonal and nonnegative
off-diagonal entries. Thus $e^{Ht}\geq0$ and $-H^{-1}\geq0$, which give
\eqref{eq:comparison_solution}, the ultimate bound and $\bm c>0$.
Since $\bm c^{\trans}H=-\bm1^{\trans}$,
$D^+V_c\leq-\bm1^{\trans}\bm\chi+\bm c^{\trans}\bm d$.
Equation~\eqref{eq:component_extraction} is the ellipsoid support in
$P\otimes I_2$ after undoing the position normalization.
\end{IEEEproof}

The core contains two ego-velocity states and six states per target. Its
certificate requires neither informative yaw measurements nor a yaw
innovation chart. The shared gyro also couples the two core stages;
\eqref{eq:cascade_comparison} uses a triangular norm bound and does not
claim to retain all cross-stage cancellations. Orientation and position
remain separately certified outputs.

\subsection{Gain Selection}
\label{subsec:gainSelection}

With $T_{\rm dom}=\rho_{\max}/(V_E|_{\max}+V_C|_{\max})$, the scalar
gain preference is
\begin{equation}
    \min_{k>0}k\quad\text{s.t.}\quad e^{-kT_{\rm dom}}\leq0.05,
    \qquad k_\psi=k_v=k_p=\frac{\log20}{T_{\rm dom}}.
    \label{eq:scalar_gain_rule}
\end{equation}
This is a transit-time preference, not an optimum of the common-error
bound in \eqref{eq:ego_rows}. Increasing $k_v$ changes the combined gyro
coefficient and need not improve that bound.

For each candidate $\omega>1\,\mathrm{s}^{-1}$, fix
$\lambda_T=1/T_{\rm dom}$ and minimize $\operatorname{tr}P$ over
$P\succeq I$ and $t_q,t_s>0$ subject to
\begin{equation}
    \begin{bmatrix}
    Z&P\bm e_3&P\bm e_3\\
    \bm e_3^{\trans}P&-t_q&0\\
    \bm e_3^{\trans}P&0&-t_s
    \end{bmatrix}\preceq0,
    \label{eq:free_metric_lmi}
\end{equation}
where
$Z=\omega(A_l^{\trans}P+PA_l)+2\lambda_TP
+\diag(0,t_qL_q^2/\omega^2,t_sL_s^2)$.
Its Schur complement is Lemma~\ref{lem:dissipation}. The normalization
fixes the homogeneous scale while leaving the metric free; the recovered
matrix is checked in the original inequality with numerical reserves.
A bracketed scalar search then minimizes the largest ultimate physical
component bound, obtained from \eqref{eq:component_extraction}, divided
by its declared domain scale. This gives a feasible local bandwidth
selection, not a global joint optimum. Neither simulation errors nor the
sensor sample period enter the design.

\subsection{Enclosure Supplied to the Controller}
\label{subsec:enclosure}

The controller of Section~\ref{sec:control} requires a transient radius
at each observation. For the continuous observer, the same comparison
system supplies it: \eqref{eq:comparison_solution} holds for all \(t\), so
\eqref{eq:component_extraction} gives

\begin{corollary}
\label{cor:enclosure}
Under the hypotheses of Theorem~\ref{thm:cascadeISS} and any
\(\bm z_0\geq\bm\chi(0)\),
\begin{equation}
    \normtwo{\bm\rho(t)-\hat{\bm\rho}(t)}\leq B_\rho(t):=c_1z_2(t),
    \quad c_1=\sqrt{(P^{-1})_{11}},
    \label{eq:relative_position_bound}
\end{equation}
and \(\limsup_{t\to\infty}B_\rho\leq c_1(-H^{-1}\bm d)_2\). An initial
metric bound follows from componentwise initial bounds \(\bm b_{T0}\) as
\(W_T(0)\leq\bigl((D\bm b_{T0})^{\trans}|P|(D\bm b_{T0})\bigr)^{1/2}\) with
\(D=\diag(1,\omega^{-1},\omega^{-2})\).
\end{corollary}

Inertial reconstruction of the relative vector costs one further yaw term,
\begin{equation}
    \normtwo{\bm p_C-\bm p_E-R(\hat\psi_E)\hat{\bm\rho}}
    \leq B_\rho+2\bar R\sin\frac{\min(B_\psi,\pi)}2 ,
    \label{eq:inertial_relative_bound}
\end{equation}
\(\bar R=\min(\rho_{\max},\normtwo{\hat{\bm\rho}})\), and the analogous
rotation bounds convert the remaining components of
\eqref{eq:component_extraction} into the componentwise ego and target
enclosures \(\bm b_E(t)\), \(\bm b_T(t)\) that the controller reads. These
certify the current estimate; they are not predictions of future observer
output, and their propagation over the control horizon is the controller's
own obligation, treated in Section~\ref{subsec:uncertainty}. The digital
runtime additionally accounts for held measurements, radar prediction and
outages in a separate finite-time enclosure. The continuous ISS theorem
does not establish exponential stability of the sampled RK4 realization
or a directed-rounding certificate for its computed bounds.

\section{Control Design}
\label{sec:control}

One convex program is solved per sample. Its decisions are the complete
input sequence, one nonnegative sampled-CLF slack, and the separating
angles of the active encounter. Collision, actuator, slew, chart and
terminal requirements are hard. The accepted plan and its geometric
certificate are retained together for the next sample.

\subsection{Prediction Model}
\label{subsec:frenetModel}

Let
\begin{equation}
    \bm x=[s,\,d,\,e_\psi,\,v_x,\,v_y,\,r]^{\trans},\qquad
    \bm u=[\delta_F,\,\beta]^{\trans},
    \label{eq:frenet_state}
\end{equation}
where $s,d$ are station and signed lateral offset, and $e_\psi$ is heading
relative to the path tangent. At curvature $\kappa$, the exact Frenet
kinematics are
\begin{equation}
    \begin{aligned}
    \dot s&=\frac{v_x\cos e_\psi-v_y\sin e_\psi}{1-\kappa d},\\
    \dot d&=v_x\sin e_\psi+v_y\cos e_\psi,\qquad
    \dot e_\psi=r-\kappa\dot s.
    \end{aligned}
    \label{eq:frenet_kinematics}
\end{equation}
Together with \eqref{eq:bicycle_model}--\eqref{eq:fiala_force}, these
define the nonlinear vector field $f_\kappa(\bm x,\bm u)$. Passive
flat-road load in still air is
\begin{equation}
    \begin{split}
    F_R(v)={}&\tfrac12\rho_{\rm air}C_dA_fv|v|\\
       &+m_Eg(q_0+q_1|v|+q_4|v|^4)\tanh(v/v_{\rm tr}),
    \end{split}
    \label{eq:road_load}
\end{equation}
with nonnegative rolling coefficients and $v_{\rm tr}>0$. The model
uses static axle loads and excludes wheel inertia, wind, grade and dynamic
load transfer.

At prescribed longitudinal speed $v_*$, the steady turn satisfies
\begin{equation}
    \begin{gathered}
    d_*=0,\qquad e_{\psi,*}=-\operatorname{atan2}(v_{y,*},v_*),\\
    r_*=\kappa\sqrt{v_*^2+v_{y,*}^2}.
    \end{gathered}
    \label{eq:cruise_trim}
\end{equation}
The three force and moment balances determine $v_{y,*},\delta_{F,*}$ and
$\beta_*$. Front wheel forces are rotated into the body frame before
solving these balances. The resulting operating point defines
\begin{equation}
    \begin{gathered}
    A^c=D_xf_\kappa(\bm x_*,\bm u_*),\qquad
    B^c=D_uf_\kappa(\bm x_*,\bm u_*),\\
    \bm c^c=f_\kappa(\bm x_*,\bm u_*)-A^c\bm x_*-B^c\bm u_*,\\
    \dot{\bm x}=A^c\bm x+B^c\bm u+\bm c^c.
    \end{gathered}
    \label{eq:scheduled_bicycle}
\end{equation}
The Fiala and road-load Jacobians retain both steering and signed-ratio
couplings. Derivative anchors lie strictly inside $|\beta|<1$; the input
constraints retain their declared endpoints.

Each stage is integrated exactly under its held input,
\begin{equation}
    \begin{bmatrix}A_j&B_j&\bm c_j\\0&I_2&0\\0&0&1\end{bmatrix}
    =\exp\!\left(T_s\begin{bmatrix}A^c_j&B^c_j&\bm c^c_j\\0&0&0\\0&0&0\end{bmatrix}\right),
    \label{eq:held_flow}
\end{equation}
so $\bm x_{j+1}=A_j\bm x_j+B_j\bm u_j+\bm c_j$ and
$\bm x_j=F_j\bm U+\bm f_j$, with
$\bm U=[\bm u_0^{\trans},\ldots,\bm u_{N-1}^{\trans}]^{\trans}$.
A straight or constant-curvature reference uses one generator. A declared
smooth profile uses an immutable phase-indexed sequence and a
constant-curvature tail. The sampled reference defect is retained, and
replanning does not change the generators of an inherited certificate.
The certified plant is this held affine model with zero process residual;
linearization alone does not establish inclusion of the nonlinear vehicle.

\subsection{Stage Constraints}
\label{subsec:hardRows}

Every held input satisfies
\begin{equation}
    \begin{gathered}
    |\delta_F|\leq\delta_{\max},\qquad
    -1\leq\beta_{\min}\leq\beta\leq\beta_{\max}\leq1,\\
    |\bm u_j-\bm u_{j-1}|\leq T_s\dot{\bm u}_{\max},
    \end{gathered}
    \label{eq:actuator_rows}
\end{equation}
where $\bm u_{-1}$ is the actual previous input and infinite rate limits
omit the corresponding rows. Pose-domain rows keep each uncertain node
inside its certified Cartesian chart. Smooth references additionally
require a bounded station-phase error and $1-\kappa(s)d>0$ throughout the
declared reference domain.

These chart requirements are distinct from physical road edges. The
current terminal construction requires a reference without physical road
boundaries. It imposes no additional speed, tire-slip or friction-circle
constraints. Consequently a feasible affine prediction does not certify
the nonlinear tire envelope away from its operating point. Such a claim
would require a justified model-residual enclosure and compatible terminal
conditions.

\subsection{Terminal Continuation}
\label{subsec:terminal}

Let $S\bm x=[d,e_\psi,v_x,v_y,r]^{\trans}$,
$\bm e=S(\bm x-\bm x_*)$, $F_e=A_{2:6,2:6}$ and $G_e=B_{2:6,:}$.
For a constant reference,
\begin{equation}
    \bm e^+=F_e\bm e+G_e(\bm u-\bm u_*)+\bm d_*,\qquad
    F_c=F_e-G_eK,
    \label{eq:terminal_error_flow}
\end{equation}
where $K$ is the discrete Riccati feedback and $\bm d_*$ is the actual
sampled trim residual, including its numerical error. With a nonsingular
modal basis $T$, write $W_m=T^{-1}$ and define
\begin{equation}
    \mathcal Z=\{\bm e\in\R^5:|W_m\bm e|\leq\bm a\},\qquad
    \bm a>0.
    \label{eq:modal_terminal_set}
\end{equation}
Complex conjugate modes are retained, so each absolute value is a complex
magnitude and each row is an ordinary second-order cone. Let
$\bar{\bm r}$ bound future tracking-coordinate measurement error and set
\begin{equation}
    C\geq|W_mF_cT|,\qquad
    \bm d_f=|W_mG_eK|\bar{\bm r}+|W_m\bm d_*|.
    \label{eq:terminal_forcing}
\end{equation}
Synthesis verifies $C\bm a+\bm d_f\leq\bm a-2\epsilon_f\bm1$,
together with actuator-amplitude and slew support inequalities.

\begin{lemma}
\label{lem:terminal}
If these inequalities hold, the hypothetical held feedback
$\bm\kappa(\hat{\bm e})=\bm u_*-K\hat{\bm e}$ keeps the true error in
$\mathcal Z$ under $\hat{\bm e}=\bm e+\bm\nu$,
$|\bm\nu|\leq\bar{\bm r}$. The previous input and admissibility of the
first feedback transition are part of the terminal information state.
\end{lemma}

\begin{IEEEproof}
The true successor is $F_c\bm e-G_eK\bm\nu+\bm d_*$, hence its modal
magnitudes are at most $C\bm a+\bm d_f$. For a freely optimized input
$\bm u=\bm\kappa(\hat{\bm e})+\bm v$, impose
\begin{equation}
    |W_mG_e\bm v|\leq\bm a-C\bm a-\bm d_f-\epsilon_f\bm1.
    \label{eq:terminal_successor_cone}
\end{equation}
Let $H_u=[I;-I]$ and let $\bm b_h$ be the actuator box shifted by trim
and reduced by its numerical reserve. The rows
\begin{equation}
    H_u\bm v\leq\bm b_h-|H_uKT|\bm a-|H_uK|\bar{\bm r}
    \label{eq:terminal_input_support}
\end{equation}
ensure amplitude admissibility. With $J_u=I+KG_e$, define
\begin{equation}
    \bm s_u=|K(I-F_c)T|\bm a
      +( |J_uK|+|K| )\bar{\bm r}+|K\bm d_*|.
    \label{eq:terminal_rate_support}
\end{equation}
The next hypothetical feedback differs from the current input by
$K(I-F_c)\bm e+J_uK\bm\nu-J_u\bm v-K\bm\nu^+-K\bm d_*$.
Thus $\pm J_u\bm v\leq T_s\dot{\bm u}_{\max}-\bm s_u$, together
with current slew, guarantees the next transition. Synthesis makes
$\bm v=0$ feasible in all these rows. Every accepted $\bm v$ preserves
true-state membership and the next feedback candidate's admissibility.
\end{IEEEproof}

Entry from the finite plan with terminal center $\bm z_N$ and radius
$\bm\rho_N$ requires
\begin{equation}
    \begin{gathered}
    |W_mS(\bm z_N-\bm x_*)|+|W_mS|\bm\rho_N
        \leq\bm a-\epsilon_f\bm1,\\
    |\bm u_*-KS(\bm z_N-\bm x_*)-\bm u_{N-1}|
        +|KS|\bm\rho_N\leq T_s\dot{\bm u}_{\max}.
    \end{gathered}
    \label{eq:terminal_entry}
\end{equation}
The invariant object includes true-state membership in $\mathcal Z$ and
the conditioned measurement set; a rectangular outer hull need not itself
be invariant. Velocity uncertainty is therefore admitted without forcing
every terminal state to rest. After encounter completion, the terminal
law supplies a feasible candidate to a one-hold optimization whose input
remains free.

For a smooth reference, the terminal family also includes station phase.
The six-dimensional indexed sets satisfy
$C_j\bm a_j+\bm d_{f,j}\leq\bm a_{j+1}-\bm\epsilon_{f,j}$, with
reference defects, cross-stage gains, slew and phase-domain supports
included. The final constant-curvature set is invariant under reference
station translation. The same induction then applies to this declared
family; arbitrary changes of reference are new obligations.

\subsection{Joint Node Separation}
\label{subsec:geometry}

Separation is imposed in Cartesian space. At each node a certified chart
gives an affine relative center and ego yaw,
\begin{equation}
    \bm r(\bm x)=\bm c_p+P_p\bm x,\qquad
    \psi(\bm x)=\psi_0+\bm w^{\trans}\bm x.
    \label{eq:affine_chart}
\end{equation}
Let $B_E,B_T$ be the convex hulls of the ego and target rectangles over
their yaw-error intervals, expressed about their nominal orientations.
Position uncertainty is enclosed by $G[-1,1]^m$ and a chart-error disk of
radius $\rho_p$. With $\bm n(\theta)=[\cos\theta,\sin\theta]^{\trans}$
and fixed target yaw center $\phi$, define
\begin{equation}
    \begin{split}
    f(\bm x,\theta)={}&d_c+\rho_p
       +h_{B_E}(\bm n(\theta-\psi(\bm x)))\\
       &+h_{B_T}(\bm n(\theta-\phi))
       +\|G^{\trans}\bm n(\theta)\|_1
       -\bm n(\theta)^{\trans}\bm r(\bm x).
    \end{split}
    \label{eq:collision_residual}
\end{equation}
Here $h_B$ denotes support. For collision records $d_c=0$: there is no
additional physical clearance buffer. Numerical reserves enforce strictly
positive separation after independent verification. Uncertainty and chart
remainders remain part of the occupied sets. The lifted-certificate
viewpoint is consistent with exact convex-body separation
\cite{ZhangLinigerBorrelli2018}; optimizing trajectory and direction
together remains nonconvex.

\begin{lemma}
\label{lem:separationRow}
At an anchor $(\bm x^0,\theta^0)$, put
$\Delta\bm r=P_p\Delta\bm x$, $\Delta\psi=\bm w^{\trans}\Delta\bm x$,
$\bm t(\theta)=[-\sin\theta,\cos\theta]^{\trans}$, and impose
$\normtwo{\Delta\bm r}\leq R_{\rm step}$.
Let $R_E,R_T$ be body circumradii,
$R_U=\sum_i\normtwo{G_i}$, $R_r=\normtwo{\bm r^0}+R_{\rm step}$,
and $\eta=1/R_{\rm step}$. Define the affine arguments
\begin{equation}
    \begin{aligned}
    \bm a_E&=\bm n(\theta^0-\psi^0)
           +\bm t(\theta^0-\psi^0)(\Delta\theta-\Delta\psi),\\
    \bm a_T&=\bm n(\theta^0-\phi)+\bm t(\theta^0-\phi)\Delta\theta,\\
    \bm a_U&=\bm n^0+\bm t^0\Delta\theta.
    \end{aligned}
    \label{eq:support_arguments}
\end{equation}
Then the convex function
\begin{equation}
    \begin{aligned}
    \widehat f^0={}&d_c+\rho_p+h_{B_E}(\bm a_E)+h_{B_T}(\bm a_T)
        +\|G^{\trans}\bm a_U\|_1\\
       &-(\bm n^0)^{\trans}\bm r^0-(\bm n^0)^{\trans}\Delta\bm r
          -(\bm t^0)^{\trans}\bm r^0\Delta\theta\\
       &+\tfrac12 R_E(\Delta\theta-\Delta\psi)^2
          +\tfrac12\eta\normtwo{\Delta\bm r}^2\\
       &+\tfrac12(R_T+R_U+R_r+\eta^{-1})(\Delta\theta)^2
    \end{aligned}
    \label{eq:collision_majorant}
\end{equation}
satisfies $f\leq\widehat f^0$ on the step domain, with equality at the
anchor.
\end{lemma}

\begin{IEEEproof}
The unit circle satisfies
$\normtwo{\bm n(\alpha+\xi)-\bm n(\alpha)-\bm t(\alpha)\xi}
\leq\xi^2/2$. A radius-$R$ body's support is $R$-Lipschitz, giving
the ego, target and zonotope remainders. The center remainder is at most
$R_r(\Delta\theta)^2/2$, and Young's inequality bounds the cross term by
$(\eta\normtwo{\Delta\bm r}^2+\eta^{-1}(\Delta\theta)^2)/2$.
All remainders vanish at the anchor.
\end{IEEEproof}

For a nondegenerate rectangle, the yaw hull is the intersection of its
circumdisk and the chord halfspaces left by its rotated vertex arcs:
$B=\{\bm y:\normtwo{\bm y}\leq R,D\bm y\leq\bm b\}$.
Strong duality gives the exact support epigraph
\begin{equation}
    h_B(\bm v)=\min_{\bm\lambda\geq0}
       \bm b^{\trans}\bm\lambda+R\normtwo{\bm v-D^{\trans}\bm\lambda}.
    \label{eq:yaw_support_dual}
\end{equation}
Zero-width yaw intervals use the rectangle's absolute-value support.
These epigraphs and the quadratic terms in \eqref{eq:collision_majorant}
have second-order cone representations.

Each target also has a terminal exit record: use a point ego body and
$d_c=R_{\rm sensing}+\epsilon_{\rm encounter}$ in
\eqref{eq:collision_residual}. The resulting halfspace places the entire
target footprint outside the sensing disk by a stored absolute deadline.
The active horizon shrinks toward this deadline. A sound current scan or
current conditioned enclosure must confirm release; prediction alone does
not remove a target, and release does not prove it can never return.

\subsection{Relaxed Sampled CLF}
\label{subsec:clf}

Let $P=R^{\trans}R\succ0$ be the discrete Riccati metric for
\eqref{eq:terminal_error_flow}, normalized by a positive scalar, and let
\begin{equation}
    V(\bm e)=\bm e^{\trans}P\bm e,\qquad
    F_c^{\trans}PF_c\preceq q_0P,\qquad q_0<1.
    \label{eq:clf_definition}
\end{equation}
For $0<\varsigma<1$, set $c=1-\varsigma(1-q_0)$ and
$a=(c+q_0)/2<c$. The first held input satisfies
\begin{equation}
    \begin{gathered}
    \hat{\bm e}^+=F_e\hat{\bm e}+G_e(\bm u_0-\bm u_*)+\bm d_*,\\
    \normtwo{R\hat{\bm e}^+}
      \leq\sqrt a\normtwo{R\hat{\bm e}}+\epsilon_{\rm num}+\delta,
    \qquad\delta\geq0.
    \end{gathered}
    \label{eq:clf_row}
\end{equation}
For $\bm e=\hat{\bm e}+\bm\eta$, $|\bm\eta|\leq\bm r_e$, put
$\sigma=\sqrt a\normtwo{|R|\bm r_e}
+\normtwo{|RF_e|\bm r_e}+2\epsilon_{\rm num}$.
The triangle inequality followed by Young's inequality gives
\begin{equation}
    \begin{gathered}
    V(\bm e^+)\leq cV(\bm e)+B+S_\delta,\\
    B=\frac{c\sigma^2}{c-a},\qquad
    S_\delta=\frac{c}{c-a}\delta(2\sigma+\delta).
    \end{gathered}
    \label{eq:sampled_dissipation}
\end{equation}
Positive slack permits $V$ to increase. Strict decrease follows when
$B+S_\delta<(1-c)V$; if a uniformly bounded allowance $\bar E$ exists
along an indefinitely feasible constant-reference sequence, then
$\limsup V\leq\bar E/(1-c)$. A positive slack penalty alone proves
neither that allowance nor asymptotic recovery. For scheduled references,
the cone uses the next reference and $P_{j+1}$ on its left and $P_j$ on
its right, with the full affine reference defect included.

\subsection{The Program}
\label{subsec:program}

Let $\mathcal D$ collect the affine dynamics, actuator and chart rows,
terminal constraints and \eqref{eq:clf_row}. With anchor
$q^0=(\bm U^0,\delta^0,\bm\theta^0)$, solve
\begin{equation}
    \begin{aligned}
    \min_{q\in\mathcal D}\quad&
      \sum_{j=1}^{N}\bm e_j^{\trans}P_j\bm e_j
      +\sum_{j=0}^{N-1}(\bm u_j-\bm u_{*,j})^{\trans}
          R_u(\bm u_j-\bm u_{*,j})\\
      &+\rho\delta^2+\tfrac{\mu}{2}\normtwo{D_q(q-q^0)}^2\\
    \text{s.t.}\quad&\widehat f_i^0(q)\leq b_i\leq0,\qquad\forall i,\\
      &\normtwo{\Delta\bm r_i}\leq R_{\rm step},\quad
        |\Delta\theta_i|\leq\Theta_{\rm step}.
    \end{aligned}
    \label{eq:controller_program}
\end{equation}
Here $R_u\succ0$, $\rho>0$, and the small proximal weight $\mu>0$
scales input and slack displacements by $D_q$, with unit angle scale.
Tracking and input effort cover the complete horizon, and input effort is
relative to the force-balanced trim. The $b_i$ contain the retained
numerical reserves. The only slack is the performance variable $\delta$.
The joint program is a convex SOCP with quadratic objective.

An existing encounter uses its shifted complete certificate as the
anchor. Fresh admission uses the available nominal plan and geometric
directions. There is one hard convex restriction, without restoration,
alternate angle initializations or a second performance solve. The
restriction may be empty even when the original nonconvex problem has a
solution. Failure therefore means that no certificate was obtained in
this restriction, not that collision is unavoidable.

The independent verifier reconstructs the trajectory from the input plan
and checks physical rows, terminal and CLF cones, and the original
nonlinear residuals \eqref{eq:collision_residual}. Solver status alone
does not authorize execution. A failed improvement may retain a
separately verified incumbent; fresh failure without an incumbent issues
no command. The complete frame must meet its acceptance deadline in
either case. One solve per frame bounds the number of optimization calls,
not their worst-case elapsed time.

\subsection{Recursive Feasibility}
\label{subsec:feasibility}

Controller state contains the accepted plan, occupied sets, angles,
charts, generators, numerical reserves, terminal witness and absolute
exit deadline. After applying $\bm u_0^\star$, the candidate is the
remaining suffix,
\begin{equation}
    \widetilde{\bm U}=[(\bm u_1^\star)^{\trans},\ldots,
        (\bm u_{N-1}^\star)^{\trans}]^{\trans},\qquad
    \widetilde{\bm\theta}=\operatorname{shift}(\bm\theta^\star).
    \label{eq:shifted_candidate}
\end{equation}
No new tail is appended while an encounter retains its exit deadline.

\begin{proposition}
\label{prop:recursive}
Assume feasible initial admission, exact execution of \eqref{eq:held_flow},
sound measurements within their admitted bounds, unchanged model,
reference and sensing contracts, and target motion inside its admitted
finite-flow sets. Suppose every new or enlarged obligation admits a new
certificate and release is currently confirmed by the stored deadline.
Then retaining the complete certificate and using its touching majorants
makes each compatible successor subproblem nonempty. Every accepted
command satisfies the certified node and held-input constraints.
\end{proposition}

\begin{IEEEproof}
Substituting the executed input in every stored row, state map and cone
leaves the accepted suffix feasible. Conditioning narrows the covered
sets, and confirmed release only removes the corresponding target's
obligations. For each remaining inherited angle, narrowing an occupied
set cannot increase its support residual. Lemma~\ref{lem:separationRow}
touches at this complete feasible suffix, so it lies in the next convex
restriction, including the zero-displacement step domain. A finite
nonnegative $\delta$ always completes its CLF cone.

After confirmed release, a restored performance horizon replaces the
inherited one only if a concrete candidate passes every constraint of the
new program. Otherwise the suffix is retained. When it is exhausted,
\eqref{eq:terminal_entry} and Lemma~\ref{lem:terminal} supply a feasible
candidate to the terminal one-hold optimization. Every accepted input
preserves the next candidate's membership and slew admissibility.
Induction gives conditional nonemptiness. Independent acceptance verifies
the stated physical constraints at every executed step.
\end{IEEEproof}

The proposition concerns physical feasible sets and their touching
restrictions in exact arithmetic. Numerical reserve transfer and
independent verification implement the acceptance condition; they do not
guarantee timely numerical completion. The certificate covers prediction
nodes only. Two rectangles may be separated at both endpoints and
intersect during the held input, as the trials below demonstrate.

\subsection{State Uncertainty}
\label{subsec:uncertainty}

The controller receives a current box
$\bm x_0\in\bm z_0+[-\bm\rho_0,\bm\rho_0]$. For its zero-residual
declared plant, propagation is
\begin{equation}
    \bm z_{j+1}=A_j\bm z_j+B_j\bm u_j+\bm c_j,\qquad
    \bm\rho_{j+1}=|A_j|\bm\rho_j+\bm a_j,
    \label{eq:box_propagation}
\end{equation}
where $\bm a_j$ is an arithmetic allowance. An affine row is tightened by
its support,
\begin{equation}
    G_j\bm z_j+H_j\bm u_j+|G_j|\bm\rho_j\leq\bm h_j.
    \label{eq:robust_row}
\end{equation}
Target tubes propagate the admitted position, velocity, acceleration,
yaw and yaw-rate sets with Cartesian jerk and yaw-acceleration bounds.
They predict finite future occupancy; a current observer radius is not
substituted for that prediction.

At the next observation, let $\bm\Delta=\hat{\bm x}-\bm z^-$ and let
$\hat{\bm x}+[-\bm b,\bm b]$ be the new measurement set. About the
retained center, compute
\begin{equation}
    \begin{gathered}
    \underline{\bm\eta}=\max(-\bm\rho^-,\bm\Delta-\bm b),\qquad
    \overline{\bm\eta}=\min(\bm\rho^-,\bm\Delta+\bm b),\\
    \bm\rho^+=\min\{\bm\rho^-,
         \max(|\underline{\bm\eta}|,|\overline{\bm\eta}|)\}.
    \end{gathered}
    \label{eq:box_intersection}
\end{equation}
A nonempty intersection contains the truth and its symmetric outer box
stays inside the inherited box. No hypothetical future measurement
contraction is used in the finite plan. An inconsistent observation or an
enlarged sensing bound requires rejection or new admission. The modal
terminal construction admits bounded velocity error, but neither it nor
the estimator's continuous theorem certifies arbitrary model mismatch,
missed detections or new unseen traffic.

\section{Results}
\label{sec:results}

\subsection{Synthesized Design}
\label{subsec:designResults}

The estimator premises are $V_E\in[5,20]$\,m/s,
$\bar\omega_E=0.30$\,rad/s; $V_C\in[10,20]$\,m/s,
$A_{\max}=2$\,m/s$^2$, $\beta_{\max}=0.015$\,rad,
$l_{r,C}=1.6$\,m, $\rho_{\max}=50$\,m; $l_{r,E}=1.45$\,m,
$\beta_{\rm dom}=0.05$\,rad, $\bar d_{\rm st}=0.02$\,rad/s; and
$\bar n_p=0.12$\,m, $\bar n_v=0.02$\,m/s,
$\bar n_a=0.05$\,m/s$^2$, $\bar n_g=0.002$\,rad/s and
$\bar n_R=0.12$\,m. Thus $\Omega_{\max}=0.1875$\,rad/s,
$S=4.2499$\,m/s$^2$ and $T_{\rm dom}=1.25$\,s.

Rule~\eqref{eq:scalar_gain_rule} gives
$k_\psi=k_v=k_p=2.3966$\,s$^{-1}$, while
Lemma~\ref{lem:lipschitz} gives $L_q=0.8011$\,s$^{-2}$ and
$L_s=0.8708$\,s$^{-1}$. The normalized injection is
$\bm l=[4.5756,7.9010,4.7805]^{\trans}$.
Freeing the metric in \eqref{eq:free_metric_lmi} reduces the selected
bandwidth from $6.4431$ to $2.5910$\,s$^{-1}$ at the same required decay
$\lambda_T=0.8$\,s$^{-1}$. The physical innovation gains become
$[l_1\omega,l_2\omega^2,l_3\omega^3]=[11.8555,53.0434,83.1565]$,
and the continuous relative-position ultimate bound becomes $2.2072$\,m,
compared with $5.7599$\,m for the preceding fixed-metric design of the
same direct-velocity observer. This bound is distinct from the online
digital enclosure and from observed tracking error.

For the free-metric design, the numerical comparison is
\begin{equation}
    H=\begin{bmatrix}-2.3966&0\\4.2972&-0.8000\end{bmatrix},\qquad
    \bm d=\begin{bmatrix}0.20066\\4.96511\end{bmatrix}.
    \label{eq:numeric_comparison}
\end{equation}
The resulting body-velocity and independent position ultimate radii are
0.08373\,m/s and 0.12835\,m. The target velocity and acceleration radii
are 12.729\,m/s and 30.263\,m/s$^2$; in particular, the latter exceeds
the target acceleration domain. These worst-case values remain too broad
to represent typical estimation accuracy.

The paired gain comparison uses identical truth, sensor draws and initial
conditions, with ten seeds (71--80), 12\,s trials and errors evaluated
after the first 2\,s. Synchronized measurements arrive at 50\,Hz, except
for the 25\,Hz case, and the maximum RK4 step is 0.005\,s. The changing
motion case adds a smooth acceleration and curvature transition at 6\,s.
Table~\ref{tab:estimatorComparison} reports means of the per-seed RMSEs;
the noise-free row is one paired trial. Both designs use the same NRMM
and digital realization, so the comparison isolates the gain synthesis.

\begin{table}[!t]
\caption{Paired estimator RMSE with fixed and free Lyapunov metrics.}
\label{tab:estimatorComparison}
\centering
\small
\begin{tabular}{llrrr}
\toprule
Case & Metric & Position & Velocity & Accel.\\
     & & (m) & (m/s) & (m/s$^2$)\\
\midrule
Noise, 50 Hz & Fixed & 0.04762 & 0.4820 & 1.828\\
             & Free  & 0.02946 & 0.1214 & 0.185\\
Changing motion & Fixed & 0.04625 & 0.4443 & 1.602\\
                & Free  & 0.02994 & 0.1325 & 0.310\\
Noise, 25 Hz & Fixed & 0.07152 & 0.7145 & 2.672\\
             & Free  & 0.04285 & 0.1772 & 0.267\\
Noise-free & Fixed & 0.000008 & 0.01506 & 0.00571\\
           & Free  & 0.000123 & 0.01531 & 0.00645\\
\bottomrule
\end{tabular}
\end{table}

At 50\,Hz the noisy velocity and acceleration RMSE fall by 74.82\% and
89.90\%, respectively. The noise-free errors increase slightly, and
changing-target response is slower. In a separate eleven-scenario
campaign with fresh seeds 101--110, the nominal relative-position RMSE
is $0.02980\pm0.00183$\,m, velocity RMSE is
$0.1231\pm0.0070$\,m/s, and acceleration RMSE is
$0.1873\pm0.0102$\,m/s$^2$, where the deviations are across ten seeds.
Curvature-response lag is approximately 0.6\,s. A one-second radar outage
produces a worst position error of 0.3315\,m and an enclosure radius up
to 87.7\,m; the five shorter 0.2\,s gaps have a worst position error of
0.1250\,m. These results quantify the bandwidth and outage tradeoffs.

All 55,701 exported samples in the 101-trial campaign lie inside their
published relative-position enclosures. This includes the deliberate
case with twice the declared sensor noise, for which empirical
containment is not a certified property because the noise premise fails.
The nominal mean enclosure is approximately 1.76\,m, well above its
typical error. The observations establish neither a sampled exponential
stability theorem nor validation on real sensors.

\subsection{Closed Loop}
\label{subsec:controlResults}

The controller records distinguish the current single-restriction policy
from its preceding joint-support admission search. The latter used
feasibility restoration and multiple initializations at admission, followed
by one hard improvement on compatible successors. Its September 19
campaigns supply the physical-gap and timing measurements below. Those
admission results do not validate the current policy's success rate: the
single restriction in \eqref{eq:controller_program} can reject an
encounter that the preceding search admitted.

Fifteen deterministic trials combine straight motion and curvatures
$\kappa=\pm0.01,\pm0.02$\,m$^{-1}$ with stationary, oncoming and
crossing targets. Each requests 300 held inputs at $T_s=0.1$\,s,
$v_{\rm ref}=8$\,m/s and sensing radius 16\,m, with seed 20260912.
The ego plant is the declared exact affine model, sensing is exact, and
process residual and target jerk are zero. No physical road boundaries
or estimator are active. Offline search and frame budgets are 30\,s;
these trials therefore study the algorithm without establishing timely
execution at the physical hold period.

With no added physical clearance buffer, all fifteen trials numerically
complete and all 4,500 commands pass their hard node certificates.
Eleven equally spaced physical checks per hold nevertheless detect
intersection in the straight stationary and oncoming cases. The other
thirteen cases retain positive sampled body gap. Table~\ref{tab:nodeGap}
reports the two failures and an independent 1,001-point audit of each
worst hold. All four rectangle edge-normal projections overlap at the
reported negative minima, confirming physical intersection.

\begin{table}[!t]
\caption{Node separation and inter-node collision without an added buffer.}
\label{tab:nodeGap}
\centering
\begin{tabular}{lrrr}
\toprule
Target & Minimum node & 11-point & Dense hold\\
       & gap (m) & gap (m) & gap (m)\\
\midrule
Stationary & $7.9880\times10^{-5}$ & $-0.22537$ & $-0.23385$\\
Oncoming & $6.9076\times10^{-5}$ & $-0.05927$ & $-0.05942$\\
\bottomrule
\end{tabular}
\end{table}

The dense minima occur at 1.5254\,s and 3.7631\,s, inside holds whose
two endpoint gaps are positive. The twelve circular trials have a
smallest observed gap of 0.01429\,m, but finite physical samples do not
prove continuous-time noncollision in those cases either. Thus numerical
completion, node certification and physical collision avoidance are
different outcomes.

A separate warmed timing campaign repeats all fifteen cases twice after
900 excluded warmup holds. It uses MATLAB R2026a on an AMD Ryzen 7
7800X3D with one computational thread and the profiler disabled. Frame
time includes input construction, the complete controller, all solves
and independent acceptance; plant propagation and physical audits are
outside that timer. The solve-phase timer includes its transcription and
solver-interface overhead. Table~\ref{tab:controllerRuntime} separates
the principal frame categories.

\begin{table}[!t]
\caption{Preceding joint-support search: warmed complete-frame time.}
\label{tab:controllerRuntime}
\centering
\begin{tabular}{lrrr}
\toprule
Frame category & Count & Median (ms) & Maximum (ms)\\
\midrule
Active continuation & 986 & 15.816 & 48.790\\
Target-free continuation & 7940 & 4.480 & 6.470\\
Admission & 32 & 655.210 & 4653.899\\
\bottomrule
\end{tabular}
\end{table}

All 986 active continuations use one conic solve and finish within
100\,ms. Their solve-phase median and maximum are 6.164 and 22.768\,ms.
By contrast, 27 of 32 admissions exceed 100\,ms. Ten initial target-free
frames and 32 release transitions complete the 9,000-frame record;
the two straight inter-node collisions recur in each repetition. These
are observed timings for the preceding search, not a worst-case execution
bound or measurements of the current single-restriction admission.

The current restriction preserves trajectory and angle decisions while
removing the repeated admission solves. Its local nature remains
essential. For relative center $(0.9,y)$, $0\leq y\leq2$, square
obstacle $[-1,1]^2$ and a 0.1\,m separation requirement, the anchor
$(y,\theta)=(0,0)$ gives a strictly positive majorant for every admissible
increment, although $(1.2,\pi/2)$ satisfies the original separation
condition. This geometric counterexample, also encoded in the regression
tests, shows why one convex restriction cannot be assigned a global
admission guarantee. The 0.1\,m requirement belongs to that illustrative
example and is not a controller clearance setting.

The available estimator-in-the-loop record is partial and predates the
joint-support revision: on the declared affine plant it completes 36
holds, then stops at 3.6\,s at the first frame after target acquisition.
It does not validate closed-loop avoidance by the current combined
algorithm. Earlier 14-degree-of-freedom plant trials use a different
controller and are not evidence for the present affine-plant certificate.

\section{Discussion}
\label{sec:discussion}

The two certificates act at different levels. Theorem~\ref{thm:cascadeISS}
is continuous-time and concerns the velocity--target core. Keeping yaw
outside that cascade removes its chart conditions from core stability,
while orientation remains necessary for inertial reconstruction. The free
metric permits lower target bandwidth at the same certified decay floor;
the paired trials show reduced derivative noise together with increased
maneuver lag. The large difference between observed errors and published
enclosures remains relevant to controller admission.

Proposition~\ref{prop:recursive} preserves a complete admitted
certificate under compatible updates. It does not imply that the initial
convex restriction is nonempty, that an unseen target can always be
admitted, or that a numerical answer arrives before the actuator deadline.
The modal terminal set supports bounded velocity uncertainty, but its
current reference domain excludes physical road boundaries and requires
the declared zero-residual affine plant.

The sampled CLF bounds the actual one-step Lyapunov change with its
slack-dependent allowance. It therefore avoids interpreting a
continuous-time derivative row as sampled dissipation. Positive slack
still permits tracking error to grow. Likewise, joint separation at the
nodes does not cover motion within a hold: the observed straight-road
intersections are consistent with the certificate's stated scope. A
whole-hold occupied-set certificate, successful admission within the
execution period, and a complete estimator--controller experiment remain
necessary before claiming physical closed-loop collision avoidance.

\section{Conclusion}
\label{sec:conclusion}

Retaining measured acceleration and yaw rate as exogenous inputs removes
the zero-jerk and zero-yaw-acceleration assumptions of the differentiated
ego design. The direct velocity measurement in \eqref{eq:ego_observer}
keeps the common gyro error in one column and yields the two-stage
comparison of Theorem~\ref{thm:cascadeISS}. A parallel yaw observer and
orientation set supply inertial reconstruction. The free-metric target
certificate retains the covariant NRMM and its analytic Lipschitz bounds;
in the paired nominal trials it reduces velocity and acceleration RMSE
by 74.82\% and 89.90\%, respectively, at the cost of a slower response.

The controller jointly optimizes the trajectory and separating angles in
\eqref{eq:controller_program}, with a soft sampled CLF and hard node,
actuator, slew, chart and terminal conditions. A fixed encounter-exit
deadline and robust modal cruise set replace infinite-future rest
separation. Retaining the complete shifted certificate gives conditional
recursive subproblem feasibility without a second optimization. Initial
admission, timely numerical completion and inter-node safety remain
separate requirements; the recorded physical audits demonstrate why the
last cannot be inferred from positive node gaps.

\section*{Data and Code Availability}

No external recorded dataset was used. All numerical results in this draft
were generated by deterministic or seeded synthetic MATLAB scripts in the
accompanying repository. A public repository URL, release commit, license,
and archival identifier should be inserted before submission.

\section*{Ethics Statement}

The reported work used only software tests and synthetic simulations. It did
not involve human participants, animals, personal data, or public-road
experiments.

\section*{Author Contributions}

CRediT roles and author-specific contributions were not available during
draft preparation and must be completed by the authors before submission.

\section*{Funding and Competing Interests}

Funding and competing-interest declarations were not available during draft
preparation. The authors must replace this placeholder with complete
statements before submission.

\section*{AI Assistance Disclosure}

This draft was prepared with assistance from a generative AI system for
repository analysis, literature organization, original prose drafting,
LaTeX formatting, and consistency checks. The quantitative values are drawn from the versioned experiment records
and checked against their numerical exports. The authors remain responsible
for verifying every equation, citation, result, interpretation, and
submission declaration.

\bibliographystyle{IEEEtran}
\bibliography{references}

\end{document}
